% ---------- δικές μας μακροεντολές για τα διαγράμματα του Γ1 ----------
\definecolor{gaC1}{HTML}{16695A}
\definecolor{gaC2}{HTML}{E8A33D}
\definecolor{gaC3}{HTML}{D1624B}
\definecolor{gaC4}{HTML}{5B8DB8}
\definecolor{gaC5}{HTML}{9A7BB5}
\definecolor{gaC6}{HTML}{8FB86B}
\definecolor{gaC7}{HTML}{C9C9C9}

% ραβδόγραμμα: \gaRabdos{μέγιστο y}{βήμα y}{ετικέτα/τιμή,...}{τίτλος άξονα y}
\newcommand{\gaRabdos}[4]{%
  \pgfmathsetmacro{\gaYs}{3.3/#1}%
  \begin{tikzpicture}[x=1.45cm,y=\gaYs cm,baseline=0cm]
    \foreach \gat in {0,#2,...,#1}{%
      \draw[gray!35,line width=0.3pt] (0.2,\gat)--(7.4,\gat);
      \node[left,font=\scriptsize] at (0.2,\gat) {\gat};}
    \draw[thick] (0.2,0)--(0.2,#1);
    \node[font=\scriptsize,rotate=90,above=0.5cm] at (0.2,{#1/2}) {#4};
    \def\gaMaxI{0}%
    \foreach \gal/\gav [count=\gai] in {#3}{%
      \xdef\gaMaxI{\gai}%
      \fill[gaC1!80] ({\gai+0.3-0.38},0) rectangle ({\gai+0.3+0.38},\gav);
      \node[above,font=\scriptsize] at ({\gai+0.3},\gav) {\gav};
      \node[below,font=\scriptsize,align=center,text width=1.3cm] at ({\gai+0.3},-0.04*#1) {\gal};}
    \draw[thick] (0.2,0)--({\gaMaxI+0.95},0);
  \end{tikzpicture}}

% ιστόγραμμα: \gaIstogramma{πρώτο άκρο}{πλάτος κλάσης}{μέγιστο y}{βήμα y}{συχνότητες,...}{άξονας x}{άξονας y}
\newcommand{\gaIstogramma}[7]{%
  \def\gaN{0}\foreach \gaf [count=\gaci] in {#5}{\xdef\gaN{\gaci}}%
  \pgfmathsetmacro{\gaXs}{min(1.5,9/\gaN)}%
  \pgfmathsetmacro{\gaYs}{3.3/#3}%
  \begin{tikzpicture}[x=\gaXs cm,y=\gaYs cm,baseline=0cm]
    \foreach \gat in {0,#4,...,#3}{%
      \draw[gray!35,line width=0.3pt] (0,\gat)--(\gaN+0.2,\gat);
      \node[left,font=\scriptsize] at (0,\gat) {\gat};}
    \foreach \gaf [count=\gai] in {#5}{%
      \fill[gaC2!75] ({\gai-1},0) rectangle (\gai,\gaf);
      \draw[gdark,line width=0.6pt] ({\gai-1},0) rectangle (\gai,\gaf);}
    \draw[thick] (0,0)--(0,#3);
    \draw[thick] (0,0)--(\gaN+0.2,0);
    \foreach \gak in {0,...,\gaN}{%
      \pgfmathparse{#1+\gak*#2}%
      \draw (\gak,0)--(\gak,-0.03*#3);
      \node[below,font=\scriptsize] at (\gak,-0.03*#3) {\pgfmathprintnumber[assume math mode=false,fixed,precision=2,use comma]{\pgfmathresult}};}
    \node[font=\scriptsize,below=0.45cm] at ({\gaN/2},0) {#6};
    \node[font=\scriptsize,rotate=90,above=0.55cm] at (0,{#3/2}) {#7};
  \end{tikzpicture}}

% κυκλικό διάγραμμα: \gaKykliko{ακτίνα cm}{ετικέτα/τιμή/δεύτερη γραμμή,...}
\newcommand{\gaKykliko}[2]{%
  \def\gaTot{0}\foreach \gal/\gav/\gae in {#2}{\pgfmathparse{\gaTot+\gav}\xdef\gaTot{\pgfmathresult}}%
  \begin{tikzpicture}[baseline=0cm]
    \xdef\gaA{90}%
    \foreach \gal/\gav/\gae [count=\gai] in {#2}{%
      \pgfmathsetmacro{\gaB}{\gaA-360*\gav/\gaTot}%
      \pgfmathsetmacro{\gaM}{(\gaA+\gaB)/2}%
      \fill[gaC\gai] (0,0)--(\gaA:#1)--(\gaA:#1) arc (\gaA:\gaB:#1)--cycle;
      \draw[white,line width=0.9pt] (0,0)--(\gaA:#1);
      \node[font=\scriptsize,align=center] at (\gaM:{#1+0.85}) {\gal\\ \gae};
      \xdef\gaA{\gaB}}%
    \draw[white,line width=0.9pt] (0,0)--(90:#1);
  \end{tikzpicture}}

\kefalaio{Διαχείριση δεδομένων}{Αριθμοί που μιλούν}{%
\begin{tabular}{@{}l@{\hspace{1em}}p{0.82\linewidth}@{}}
1.1 & Απογραφή, πληθυσμός, κατηγορικά και ποσοτικά δεδομένα\\[2pt]
1.2 & Συλλογή συνεχών ποσοτικών δεδομένων\\[2pt]
1.3 & Κυκλικά διαγράμματα και ιστογράμματα συχνοτήτων\\[2pt]
 & Επανάληψη του κεφαλαίου\\
\end{tabular}}
\paragrafos{Απογραφή, πληθυσμός, κατηγορικά και ποσοτικά δεδομένα}{Γ.1.1, «Απογραφή – Πληθυσμός – Κατηγορικά δεδομένα – Ποσοτικά δεδομένα»}

\begin{anagke}
Η κυρία Αλεξάνδρα ανακοινώνει ότι το σχολείο θα οργανώσει αθλητική ημέρα, αλλά το κονδύλι φτάνει για ένα μόνο άθλημα. «Ποιο να διαλέξουμε;» ρωτά. Ο Νίκος φωνάζει αμέσως «μπάσκετ!». Η Ζωή απαντά: «Εσύ το θέλεις, αλλά το σχολείο;». Ο Στέλιος προτείνει να ρωτήσουν μόνο τους φίλους του, γιατί «είναι πιο γρήγορο». Η Ρήνη αντιδρά: «Οι φίλοι σου παίζουν όλοι ποδόσφαιρο, δεν θα βγει νόημα».

Η Φανή πιάνει το τετράδιό της και προτείνει: «Να ρωτήσουμε τους συμμαθητές μας, να σημειώσουμε τις απαντήσεις και να δούμε τι βγαίνει». Αμέσως προκύπτουν ερωτήματα: ποιους θα ρωτήσουμε, τι ακριβώς θα ρωτήσουμε, πώς θα γράψουμε τις απαντήσεις ώστε να διαβάζονται με μια ματιά.

Σε αυτή την παράγραφο μαθαίνουμε να ονομάζουμε αυτό που ερευνούμε (πληθυσμός, μεταβλητή), να διαλέγουμε σωστά ποιους θα ρωτήσουμε (δείγμα), να ξεχωρίζουμε τα είδη των δεδομένων και να τα οργανώνουμε σε πίνακα συχνοτήτων και ραβδόγραμμα.
\end{anagke}

\begin{lexi}{στατιστική, απογραφή, πληθυσμός, δείγμα, μεταβλητή}{ονόματα της στατιστικής}
Η \lexiword{στατιστική} πήρε το όνομά της από το «κράτος»: οι πρώτοι που μάζευαν δεδομένα ήταν οι κυβερνήσεις, για να ξέρουν πόσοι είναι οι πολίτες τους. Η \lexiword{απογραφή} είναι «το γράψιμο από την αρχή ως το τέλος» όλων. Ο \lexiword{πληθυσμός} είναι ο «πολύς κόσμος» που μας ενδιαφέρει. Το \lexiword{δείγμα} είναι το «παράδειγμα»: ένα μικρό κομμάτι που δείχνει το όλο, όπως η γεύση ενός κουταλιού από το φαγητό. Η \lexiword{μεταβλητή} είναι αυτό που αλλάζει από άτομο σε άτομο.
\end{lexi}

\begin{metafrasi}
Όλοι όσοι μας ενδιαφέρουν. & πληθυσμός\\
Ρωτάω όλους, έναν έναν. & απογραφή\\
Ρωτάω μόνο μερικούς, επιλεγμένους με τύχη. & δείγμα\\
Αυτό που ρωτάω (ύψος, χρώμα, άθλημα). & μεταβλητή\\
Οι απαντήσεις που κατέγραψα. & στατιστικά δεδομένα\\
Πόσες φορές εμφανίστηκε μια τιμή. & συχνότητα\\
Ποιο μέρος του συνόλου είναι μια τιμή. & σχετική συχνότητα\\
\end{metafrasi}

\tmima{Η θεωρία}

\begin{orismos}{πληθυσμός και μεταβλητή}
\textbf{Πληθυσμός} είναι μια ομάδα ανθρώπων ή ένα σύνολο αντικειμένων, των οποίων τα στοιχεία μελετάμε ως προς κάποιο χαρακτηριστικό τους. Αυτό το χαρακτηριστικό ονομάζεται \textbf{μεταβλητή}.
\end{orismos}

Παραδείγματα πληθυσμών: οι μαθητές των γυμνασίων της χώρας, τα αυτοκίνητα που κινούνται σε μια εθνική οδό, τα πρόβατα του παππού Γιάννη. Παραδείγματα μεταβλητών: το ύψος των μαθητών, το χρώμα ενός αυτοκινήτου, το βάρος ενός προβάτου.

\begin{orismos}{απογραφή, δείγμα, δειγματοληψία}
Η διαδικασία με την οποία εξετάζουμε και καταγράφουμε \emph{όλα} τα στοιχεία ενός πληθυσμού λέγεται \textbf{απογραφή}. Όταν η απογραφή κοστίζει πολύ σε χρόνο ή σε χρήματα, εξετάζουμε μόνο ένα μέρος του πληθυσμού, που λέγεται \textbf{δείγμα}. Η διαδικασία επιλογής του δείγματος λέγεται \textbf{δειγματοληψία}.
\end{orismos}

Το δείγμα πρέπει να επιλέγεται έτσι ώστε τα συμπεράσματα που θα βγάλουμε να είναι \textbf{αξιόπιστα}. Τότε λέμε ότι το δείγμα είναι \textbf{αντιπροσωπευτικό} του πληθυσμού: μοιάζει με τον πληθυσμό, σαν μικρογραφία του. Αν ρωτήσουμε τον πληθυσμό ή το δείγμα, οι απαντήσεις που συγκεντρώνουμε λέγονται \textbf{στατιστικά δεδομένα} ή \textbf{παρατηρήσεις}.

\begin{kanonas}{τα είδη των δεδομένων}
\textbf{Κατηγορικά} είναι τα δεδομένα που δεν είναι αριθμοί (χρώμα ματιών, άθλημα, μάρκα κινητού).\\
\textbf{Ποσοτικά} είναι τα δεδομένα που είναι αριθμοί (ύψος, ηλικία, αριθμός αδελφών). Τα ποσοτικά χωρίζονται σε δύο είδη:\\
\textbf{διακριτά}, όταν παίρνουν μεμονωμένες τιμές (π.χ. αριθμός αδελφών: $0$, $1$, $2$, $3$ και όχι $1{,}5$)\\
\textbf{συνεχή}, όταν μπορούν να πάρουν οποιαδήποτε τιμή μέσα σε ένα διάστημα (π.χ. ο χρόνος σε λεπτά που χρειάζεται ένας μαθητής για να λύσει ένα διαγώνισμα, από $0$ ως $120$).
\end{kanonas}

\begin{center}
\begin{tikzpicture}[font=\small,
  nd/.style={draw,rounded corners=3pt,align=center,inner sep=4pt,fill=gblueL},
  nl/.style={draw,rounded corners=3pt,align=center,inner sep=4pt,fill=gwarm}]
\node[nd] (a) at (0,0) {Στατιστικά\\ δεδομένα};
\node[nd] (b) at (4,1.1) {Ποσοτικά\\ (είναι αριθμοί)};
\node[nd] (c) at (4,-1.3) {Κατηγορικά\\ (δεν είναι αριθμοί)};
\node[nl] (d) at (8.3,1.9) {Διακριτά\\ (μεμονωμένες τιμές)};
\node[nl] (e) at (8.3,0.3) {Συνεχή\\ (κάθε τιμή σε διάστημα)};
\draw (a.east)--++(0.6,0)|-(b.west);
\draw (a.east)--++(0.6,0)|-(c.west);
\draw (b.east)--++(0.6,0)|-(d.west);
\draw (b.east)--++(0.6,0)|-(e.west);
\end{tikzpicture}
\end{center}

\begin{prosoxi}
\begin{itemize}[leftmargin=1.2em,itemsep=0.3ex,topsep=2pt]
\item Δεν αρκεί να είναι κάτι γραμμένο με αριθμό για να είναι ποσοτικό. Ο αριθμός της φανέλας ενός παίκτη δεν μετρά τίποτα, απλώς τον ξεχωρίζει: είναι κατηγορικό δεδομένο. Ποσοτικό είναι ό,τι μετράει μια ποσότητα, δηλαδή απαντά στο «πόσο» ή στο «πόσα».
\item Αν το δείγμα έχει επιλεγεί από ένα μόνο τμήμα του πληθυσμού (π.χ. μόνο οι φίλοι μου, μόνο η ομάδα μπάσκετ), τα συμπεράσματα δεν είναι αξιόπιστα για όλο τον πληθυσμό.
\item Το πλήθος (αριθμός παιδιών, αριθμός αδελφών) είναι διακριτό. Το ύψος, το βάρος και ο χρόνος είναι συνεχή.
\end{itemize}
\end{prosoxi}

\begin{orismos}{συχνότητα και σχετική συχνότητα}
\textbf{Συχνότητα} μιας τιμής είναι το πλήθος των φορών που εμφανίζεται στα δεδομένα. \textbf{Σχετική συχνότητα} μιας τιμής είναι το πηλίκο της συχνότητάς της διά του συνόλου των παρατηρήσεων. Τη γράφουμε συνήθως ως ποσοστό επί τοις εκατό ($\%$).
\end{orismos}

Οι συχνότητες και οι σχετικές συχνότητες συγκεντρώνονται σε έναν \textbf{πίνακα συχνοτήτων}. Το άθροισμα όλων των συχνοτήτων είναι το σύνολο των παρατηρήσεων και το άθροισμα όλων των σχετικών συχνοτήτων είναι πάντα $100\%$ (ή $1$).

Το \textbf{ραβδόγραμμα} παριστάνει τις συχνότητες με ράβδους ίδιου πλάτους. Το ύψος κάθε ράβδου ισούται με τη συχνότητα της τιμής και οι ράβδοι δεν ακουμπούν μεταξύ τους. Είναι κατάλληλο για κατηγορικά και για διακριτά δεδομένα.

\tmima{Λυμένα παραδείγματα}

\paradeigma{Ποιος είναι ο πληθυσμός;}
\tixero Η Ρήνη θέλει να μάθει πόσα λεπτά την ημέρα παίζουν παιχνίδια στο κινητό οι $280$ μαθητές του σχολείου. Ρωτά τυχαία $30$ μαθητές από όλες τις τάξεις.\par
\tithelo Πληθυσμός, δείγμα, μεταβλητή και είδος δεδομένων.\par
\possindeo Εξετάζουμε όλους τους μαθητές του σχολείου: αυτοί είναι ο πληθυσμός. Ρωτήθηκαν μόνο $30$: αυτοί είναι το δείγμα. Αυτό που ρωτάμε είναι ο χρόνος παιχνιδιού.\par
\apantisi Πληθυσμός: οι $280$ μαθητές. Δείγμα: οι $30$ μαθητές. Μεταβλητή: τα λεπτά παιχνιδιού την ημέρα. Τα δεδομένα είναι ποσοτικά και συνεχή, γιατί ο χρόνος μπορεί να πάρει οποιαδήποτε τιμή (π.χ. $42{,}5$ λεπτά). Το δείγμα είναι αντιπροσωπευτικό, αφού τα παιδιά επελέγησαν τυχαία από όλες τις τάξεις.\par
\elegxos Το δείγμα ($30$) είναι μικρότερο από τον πληθυσμό ($280$), όπως πρέπει.

\paradeigma{Αντιπροσωπευτικό ή όχι}
\tixero Ο Ηλίας θέλει να μάθει το ύψος των μαθητών του σχολείου και μετρά τα $22$ παιδιά της ομάδας μπάσκετ της Γ΄ τάξης.\par
\tithelo Αν το δείγμα είναι αντιπροσωπευτικό.\par
\possindeo Οι παίκτες του μπάσκετ είναι συνήθως ψηλότεροι. Επιπλέον είναι μόνο της Γ΄ τάξης, που έχει μεγαλύτερα παιδιά από την Α΄ και τη Β΄.\par
\apantisi Δεν είναι αντιπροσωπευτικό. Τα συμπεράσματα δεν θα είναι αξιόπιστα για όλο το σχολείο.\par
\elegxos Αν είχε επιλέξει $22$ παιδιά με κλήρωση από όλες τις τάξεις, το δείγμα θα ήταν καλύτερο.

\paradeigma{Είδος δεδομένων}
\tixero Οι μεταβλητές: χρώμα ματιών, αριθμός αδελφών, βάρος σχολικής τσάντας.\par
\tithelo Το είδος των δεδομένων της καθεμιάς.\par
\possindeo Το χρώμα δεν είναι αριθμός. Ο αριθμός αδελφών είναι αριθμός με μεμονωμένες τιμές. Το βάρος είναι αριθμός και μπορεί να πάρει οποιαδήποτε τιμή (π.χ. $3{,}27$ kg).\par
\apantisi Χρώμα ματιών: κατηγορικά. Αριθμός αδελφών: ποσοτικά διακριτά. Βάρος τσάντας: ποσοτικά συνεχή.\par
\elegxos Το «πόσα» δίνει διακριτά, το «πόσο» δίνει συνεχή.

\paradeigma{Πίνακας συχνοτήτων και ραβδόγραμμα}
\tixero Η Φανή ρώτησε και τα $25$ παιδιά της Α$3$ ποιο είναι το αγαπημένο τους άθλημα. Ποδόσφαιρο είπαν $8$, μπάσκετ $6$, βόλεϊ $4$, κολύμβηση $3$, τένις $2$ και άλλο $2$.\par
\tithelo Τον πίνακα συχνοτήτων με τις σχετικές συχνότητες και το ραβδόγραμμα.\par
\possindeo Κάθε σχετική συχνότητα είναι η συχνότητα διά $25$: $\frac{8}{25}=0{,}32=32\%$, $\frac{6}{25}=0{,}24=24\%$, και ούτω καθεξής.\par
\apantisi

\begin{center}
\renewcommand{\arraystretch}{1.25}
\begin{tabular}{lcc}
\toprule
\textbf{Άθλημα} & \textbf{Συχνότητα} & \textbf{Σχετική συχνότητα}\\
\midrule
Ποδόσφαιρο & $8$ & $32\%$\\
Μπάσκετ & $6$ & $24\%$\\
Βόλεϊ & $4$ & $16\%$\\
Κολύμβηση & $3$ & $12\%$\\
Τένις & $2$ & $8\%$\\
Άλλο & $2$ & $8\%$\\
\midrule
\textbf{Σύνολο} & $25$ & $100\%$\\
\bottomrule
\end{tabular}
\end{center}

\begin{center}
\gaRabdos{10}{2}{Ποδόσφαιρο/8,Μπάσκετ/6,Βόλεϊ/4,Κολύμβηση/3,Τένις/2,Άλλο/2}{παιδιά}
\end{center}

\elegxos $8+6+4+3+2+2=25$ και $32+24+16+12+8+8=100$.

\tmima{Ασκήσεις}

\zoni{Το κατάλαβα}{ονόματα και είδη}

\ex Κατηγορικά ή ποσοτικά δεδομένα; Στα ποσοτικά πες αν είναι διακριτά ή συνεχή. α) Το χρώμα της φανέλας μιας ομάδας. β) Ο αριθμός των αδελφών ενός παιδιού. γ) Το αγαπημένο γλυκό. δ) Οι ώρες ύπνου μιας νύχτας. ε) Το βάρος μιας σχολικής τσάντας. στ) Ο αριθμός των επαφών στο κινητό. ζ) Η μάρκα ενός κινητού.

\ex Η Ρήνη θέλει να μάθει πόσα βιβλία διάβασαν το καλοκαίρι οι $120$ μαθητές της Α΄ τάξης του σχολείου της. Ρωτά $20$ παιδιά που διάλεξε με κλήρωση. Ποιος είναι ο πληθυσμός, ποιο το δείγμα, ποια η μεταβλητή και τι είδους είναι τα δεδομένα;

\ex Απογραφή ή δείγμα; α) Η κυρία Αλεξάνδρα καταγράφει τις απουσίες και των $25$ παιδιών της Α$3$. β) Ο δήμος ρωτά $500$ κατοίκους από τους $12\,000$ για το νέο πάρκο. γ) Ο παππούς Γιάννης δοκιμάζει ένα σταφύλι για να κρίνει αν ο αμπελώνας είναι γλυκός. δ) Ο Ηλίας μετρά το ύψος όλων των παιδιών της τάξης του.

\ex Η Αγγελική κατέγραψε τι έφεραν για κολατσιό τα $20$ παιδιά μιας τάξης. Τοστ έφεραν $7$, φρούτο $5$, κουλούρι $4$, γιαούρτι $2$ και τίποτα $2$ παιδιά. Φτιάξε πίνακα συχνοτήτων με σύνολο και σχετικές συχνότητες.

\ex Στο ραβδόγραμμα φαίνεται πώς πηγαίνουν στο σχολείο τα $25$ παιδιά της Α$3$.
\begin{center}
\gaRabdos{10}{2}{Με τα πόδια/9,Ποδήλατο/3,Λεωφορείο/5,Αυτοκίνητο/8}{παιδιά}
\end{center}
α) Ποιος τρόπος είναι ο πιο συνηθισμένος και ποιος ο πιο σπάνιος; β) Πόσα παιδιά δεν έρχονται με αυτοκίνητο; γ) Ποια είναι η σχετική συχνότητα του ποδηλάτου;

\ex Η γιαγιά Μυρτώ έχει στο χωριό $12$ κότες, $8$ πρόβατα, $6$ γίδες, $4$ γάτες και $2$ σκύλους. Φτιάξε πίνακα συχνοτήτων και ραβδόγραμμα. Πόσα ζώα έχει συνολικά;

\ex Σωστό ή λάθος; α) Ο πληθυσμός έχει πάντα ανθρώπους. β) Το δείγμα έχει συνήθως λιγότερα στοιχεία από τον πληθυσμό. γ) Ο αριθμός των αδελφών είναι διακριτό δεδομένο. δ) Το χρώμα ματιών είναι ποσοτικό δεδομένο. ε) Στο ραβδόγραμμα το ύψος της ράβδου δείχνει τη συχνότητα.

\zoni{Το χρησιμοποιώ}{ερωτηματολόγια και πίνακες}

\ex \textbf{Το δείγμα.} Ο Στέλιος θέλει να μάθει ποιο τραγούδι προτιμούν οι μαθητές του σχολείου. Προτείνει τρεις τρόπους: α) να ρωτήσει τους $10$ φίλους του, β) να ρωτήσει $30$ παιδιά που θα διαλέξει με κλήρωση από όλες τις τάξεις, γ) να ρωτήσει όσους βρει στο κυλικείο το πρώτο διάλειμμα. Ποιος τρόπος δίνει το πιο αντιπροσωπευτικό δείγμα και γιατί;

\ex \textbf{Τα αυτοκίνητα.} Η Φανή κάθεται έξω από το σχολείο και καταγράφει το χρώμα $20$ αυτοκινήτων που περνούν: Λευκό (Α), Γκρι (Γ), Μαύρο (Μ), Κόκκινο (Κ), Άλλο (Ο).\par
\noindent Γ Κ Ο Μ Γ Μ Α Α Γ Α Μ Α Α Γ Κ Α Μ Α Α Γ\par
Φτιάξε πίνακα συχνοτήτων με σχετικές συχνότητες και ραβδόγραμμα.

\ex \textbf{Τα βιβλία.} Στο ραβδόγραμμα φαίνεται πόσα βιβλία διάβασαν το καλοκαίρι τα $30$ παιδιά μιας τάξης.
\begin{center}
\gaRabdos{10}{2}{0 βιβλία/4,1 βιβλίο/9,2 βιβλία/8,3 βιβλία/6,4 βιβλία/3}{παιδιά}
\end{center}
α) Ποια τιμή έχει τη μεγαλύτερη συχνότητα; β) Πόσα παιδιά διάβασαν τουλάχιστον $2$ βιβλία; γ) Ποια είναι η σχετική συχνότητα των παιδιών που διάβασαν ένα βιβλίο; δ) Πόσα βιβλία διάβασαν συνολικά όλα τα παιδιά;

\ex \textbf{Τα κινητά.} Ο πίνακας δείχνει τέσσερα κινητά ενός καταστήματος.
\begin{center}
\renewcommand{\arraystretch}{1.2}
\begin{tabular}{lcccc}
\toprule
\textbf{Μοντέλο} & \textbf{Χρώμα} & \textbf{Μνήμη (GB)} & \textbf{Βάρος (g)} & \textbf{Αδιάβροχο}\\
\midrule
Κ1 & μπλε & $64$ & $172{,}5$ & ναι\\
Κ2 & μαύρο & $128$ & $180$ & ναι\\
Κ3 & λευκό & $64$ & $165{,}2$ & όχι\\
Κ4 & μπλε & $256$ & $190{,}8$ & ναι\\
\bottomrule
\end{tabular}
\end{center}
Για καθεμιά από τις στήλες από τη δεύτερη ως την πέμπτη, πες αν τα δεδομένα είναι κατηγορικά ή ποσοτικά και, αν είναι ποσοτικά, διακριτά ή συνεχή.

\ex \textbf{Ο συμπληρωμένος πίνακας.} Στο ραβδόγραμμα της Ρήνης φαίνονται οι δραστηριότητες που παρακολουθούν μετά το σχολείο $40$ παιδιά. Οι ράβδοι της μουσικής, του χορού, της ζωγραφικής και της ρομποτικής έχουν ύψη $10$, $8$, $6$ και $12$. Η πέμπτη ράβδος είναι του σκακιού και δεν φαίνεται στο χαρτί. Πόσο είναι το ύψος της και ποια η σχετική συχνότητα κάθε δραστηριότητας;

\zoni{Το σκέφτομαι}{γιατί και πώς}

\ex Ο Βασίλης λέει: «Από την έρευνα της Φανής, $8$ στα $25$ παιδιά της Α$3$ προτιμούν το ποδόσφαιρο, άρα το $32\%$ όλων των μαθητών του σχολείου προτιμά το ποδόσφαιρο». Συμφωνείς; Τι θα έπρεπε να ισχύει για το δείγμα;

\ex Η Ζωή θέλει να μάθει πόσες ώρες κοιμούνται τα παιδιά της ηλικίας της. Όρισε: πληθυσμό, δείγμα, μεταβλητή, είδος δεδομένων, και πρότεινε πώς θα επιλεγεί το δείγμα.

\ex Η Αγγελική βρήκε τις σχετικές συχνότητες $40\%$, $35\%$ και $30\%$ για τρεις κατηγορίες. Ο Ηλίας λέει ότι κάτι δεν πάει καλά. Ποιος έχει δίκιο και γιατί;

\ex Σκέψου μια ερώτηση που μπορείς να κάνεις στην Α$3$ και θα δώσει: α) κατηγορικά δεδομένα, β) διακριτά ποσοτικά δεδομένα, γ) συνεχή ποσοτικά δεδομένα.

\zoni{Για το σπίτι}{ερευνώ στην οικογένεια}

\ex Ρώτησε $10$ άτομα της οικογένειας ή της γειτονιάς ποιο είναι το αγαπημένο τους φαγητό. Φτιάξε πίνακα συχνοτήτων και ραβδόγραμμα.

\ex Βρες σε εφημερίδα, περιοδικό ή ιστοσελίδα ένα διάγραμμα. Γράψε ποιος είναι ο πληθυσμός, ποια η μεταβλητή και τι είδους είναι τα δεδομένα.

\paragrafos{Συλλογή συνεχών ποσοτικών δεδομένων}{Γ.1.2, «Συλλογή συνεχών ποσοτικών δεδομένων»}

\begin{anagke}
Ο Ηλίας φέρνει στην τάξη μια μεζούρα και μετρά το ύψος και των $25$ παιδιών της Α$3$, σε εκατοστά. Γράφει τις μετρήσεις στο τετράδιό του, με τη σειρά που στάθηκαν τα παιδιά στην ουρά:

\noindent\mbox{}\hfill 160 \ 165 \ 154 \ 169 \ 171 \ 162 \ 155 \ 174 \ 157 \ 154 \ 161 \ 148 \ 163\hfill\mbox{}\\
\mbox{}\hfill 156 \ 160 \ 167 \ 158 \ 152 \ 163 \ 151 \ 150 \ 166 \ 159 \ 153 \ 158\hfill\mbox{}

«Αυτό δεν διαβάζεται!» λέει η Ρήνη. «Σχεδόν κάθε παιδί έχει διαφορετικό ύψος. Αν φτιάξω πίνακα με μια γραμμή για κάθε τιμή, θα έχω $21$ γραμμές με συχνότητα $1$ ή $2$». Η Φανή ρωτά: «Δεν μπορούμε να βάλουμε μαζί τα παιδιά με ύψος κοντά το ένα στο άλλο;».

Σε αυτή την παράγραφο μαθαίνουμε να ομαδοποιούμε τα συνεχή δεδομένα σε κλάσεις και να φτιάχνουμε πίνακα συχνοτήτων και σχετικών συχνοτήτων που διαβάζεται με μια ματιά.
\end{anagke}

\begin{lexi}{κλάση, ομαδοποίηση, κέντρο, πλάτος}{ονόματα για τις ομάδες}
Η \lexiword{κλάση} σημαίνει την «τάξη», την ομάδα: όπως η τάξη του σχολείου μαζεύει παιδιά ίδιας ηλικίας, η κλάση μαζεύει τιμές που είναι κοντά. Η \lexiword{ομαδοποίηση} είναι το «βάζω σε ομάδες». Το \lexiword{κέντρο} της κλάσης είναι το μέσο της και το \lexiword{πλάτος} της είναι πόσο φαρδιά είναι.
\end{lexi}

\begin{metafrasi}
Βάζω σε ομάδες τιμές που είναι κοντά. & ομαδοποίηση, κλάσεις\\
Η ομάδα από $150$ μέχρι $155$ (χωρίς το $155$). & κλάση $150-155$\\
Πόσο φαρδιά είναι η κλάση. & πλάτος κλάσης\\
Το μέσο της κλάσης. & κέντρο κλάσης\\
Πόσα στοιχεία έχει η κλάση. & συχνότητα της κλάσης\\
Ποιο μέρος του συνόλου είναι η κλάση. & σχετική συχνότητα της κλάσης\\
\end{metafrasi}

\tmima{Η θεωρία}

Στα συνεχή ποσοτικά δεδομένα (ύψος, βάρος, χρόνος) οι διαφορετικές τιμές είναι πάρα πολλές και σχεδόν δεν επαναλαμβάνονται. Τότε είναι προτιμότερο να χωρίσουμε τις παρατηρήσεις σε ομάδες που λέγονται \textbf{κλάσεις}. Η διαδικασία αυτή ονομάζεται \textbf{ομαδοποίηση παρατηρήσεων}.

\begin{orismos}{κλάσεις ίσου πλάτους}
Κάθε κλάση περιγράφεται από τα δύο άκρα της, π.χ. $150-155$. \textbf{Πλάτος} της κλάσης είναι η διαφορά του μικρότερου άκρου από το μεγαλύτερο: $155-150=5$. \textbf{Κέντρο} της κλάσης είναι ο αριθμός που βρίσκεται στο μέσο της: $\frac{150+155}{2}=152{,}5$. Οι κλάσεις που φτιάχνουμε έχουν \emph{ίσο πλάτος}, δεν επικαλύπτονται και καλύπτουν όλες τις παρατηρήσεις.
\end{orismos}

\begin{kanonas}{πού μπαίνει μια τιμή}
Μια παρατήρηση που συμπίπτει με το \emph{δεξί} άκρο μιας κλάσης τοποθετείται στην \emph{επόμενη} κλάση. Για παράδειγμα, η τιμή $155$ ανήκει στην κλάση $155-160$ και όχι στην $150-155$.
\end{kanonas}

\begin{kanonas}{πώς φτιάχνουμε τις κλάσεις}
$1.$ Βρίσκουμε τη μικρότερη και τη μεγαλύτερη τιμή. $2.$ Διαλέγουμε ένα πλάτος που δίνει περίπου $5$ έως $7$ κλάσεις και αριθμούς εύκολους (π.χ. $5$ ή $10$). $3.$ Ξεκινάμε από έναν στρογγυλό αριθμό λίγο μικρότερο ή ίσο με τη μικρότερη τιμή και προχωρούμε με το πλάτος ώσπου να καλύψουμε και τη μεγαλύτερη. $4.$ Μετράμε πόσες τιμές πέφτουν σε κάθε κλάση.
\end{kanonas}

\begin{kanonas}{σχετική συχνότητα}
$$\text{σχετική συχνότητα της κλάσης}=\frac{\text{συχνότητα της κλάσης}}{\text{σύνολο παρατηρήσεων}}$$
Το αποτέλεσμα γράφεται ως δεκαδικός ή, συνήθως, ως ποσοστό: $\frac{6}{25}=0{,}24=24\%$.
\end{kanonas}

Ο πίνακας με τις κλάσεις, τις συχνότητες και τις σχετικές συχνότητες λέγεται \textbf{πίνακας κατανομής συχνοτήτων}. Στην τελευταία γραμμή γράφουμε το σύνολο: το άθροισμα των συχνοτήτων είναι το πλήθος των παρατηρήσεων και το άθροισμα των σχετικών συχνοτήτων είναι $100\%$.

\begin{prosoxi}
\begin{itemize}[leftmargin=1.2em,itemsep=0.3ex,topsep=2pt]
\item Η συχνότητα και η σχετική συχνότητα ενός πίνακα με κλάσεις αναφέρονται \emph{στην κλάση} και όχι σε συγκεκριμένους αριθμούς. Αν η κλάση $165-170$ έχει συχνότητα $4$, δεν ξέρουμε πόσα παιδιά έχουν ύψος ακριβώς $167$.
\item Το δεξί άκρο δεν ανήκει στην κλάση. Το $160$ μπαίνει στην $160-165$.
\item Οι κλάσεις πρέπει να έχουν το ίδιο πλάτος και να μην αφήνουν κενά ούτε να επικαλύπτονται.
\item Το άθροισμα των σχετικών συχνοτήτων πρέπει να βγαίνει $100\%$. Αν δεν βγαίνει, έχει γίνει λάθος στην καταμέτρηση ή στη διαίρεση.
\end{itemize}
\end{prosoxi}

\tmima{Λυμένα παραδείγματα}

\paradeigma{Τα ύψη της Α$3$}
\tixero Οι $25$ μετρήσεις του Ηλία. Η μικρότερη είναι $148$ και η μεγαλύτερη $174$.\par
\tithelo Πίνακας κατανομής συχνοτήτων με κλάσεις πλάτους $5$ cm.\par
\possindeo Οι τιμές είναι από $148$ έως $174$, άρα οι κλάσεις ξεκινούν από το $145$ και φτάνουν ως το $175$: έξι κλάσεις. Μετράμε κάθε κλάση. Η $150-155$ περιέχει τα $150$, $151$, $152$, $153$, $154$, $154$: έξι τιμές. Το $155$ δεν είναι εδώ, ανήκει στην επόμενη κλάση.\par
\apantisi

\begin{center}
\renewcommand{\arraystretch}{1.25}
\begin{tabular}{cccc}
\toprule
\textbf{Κλάσεις} & \textbf{Κέντρο} & \textbf{Συχνότητα} & \textbf{Σχετική συχνότητα}\\
\midrule
$145-150$ & $147{,}5$ & $1$ & $4\%$\\
$150-155$ & $152{,}5$ & $6$ & $24\%$\\
$155-160$ & $157{,}5$ & $6$ & $24\%$\\
$160-165$ & $162{,}5$ & $6$ & $24\%$\\
$165-170$ & $167{,}5$ & $4$ & $16\%$\\
$170-175$ & $172{,}5$ & $2$ & $8\%$\\
\midrule
\textbf{Σύνολο} & & $25$ & $100\%$\\
\bottomrule
\end{tabular}
\end{center}

\elegxos $1+6+6+6+4+2=25$ και $4+24+24+24+16+8=100$.

\paradeigma{Ο χρόνος διαδρομής}
\tixero Η Φανή ρώτησε $20$ παιδιά πόσα λεπτά κάνουν για να φτάσουν στο σχολείο:\par
\noindent\mbox{}\hfill 15 \ 8 \ 14 \ 25 \ 10 \ 18 \ 11 \ 16 \ 12 \ 9 \ 7 \ 5 \ 10 \ 3 \ 6 \ 20 \ 13 \ 4 \ 22 \ 9\hfill\mbox{}\par
\tithelo Πίνακας με $6$ κλάσεις ίσου πλάτους.\par
\possindeo Οι τιμές είναι από $3$ έως $25$. Διαλέγουμε πλάτος $5$ και ξεκινάμε από το $0$: $0-5$, $5-10$, $10-15$, $15-20$, $20-25$, $25-30$. Η τιμή $5$ πάει στην $5-10$, οι $10$ στην $10-15$, η $15$ στην $15-20$, η $20$ στην $20-25$ και η $25$ στην $25-30$. Σε κάθε σύνολο μετράμε: $\{3,4\}$, $\{5,6,7,8,9,9\}$, $\{10,10,11,12,13,14\}$, $\{15,16,18\}$, $\{20,22\}$, $\{25\}$.\par
\apantisi

\begin{center}
\renewcommand{\arraystretch}{1.25}
\begin{tabular}{cccc}
\toprule
\textbf{Κλάσεις (λεπτά)} & \textbf{Κέντρο} & \textbf{Συχνότητα} & \textbf{Σχετική συχνότητα}\\
\midrule
$0-5$ & $2{,}5$ & $2$ & $10\%$\\
$5-10$ & $7{,}5$ & $6$ & $30\%$\\
$10-15$ & $12{,}5$ & $6$ & $30\%$\\
$15-20$ & $17{,}5$ & $3$ & $15\%$\\
$20-25$ & $22{,}5$ & $2$ & $10\%$\\
$25-30$ & $27{,}5$ & $1$ & $5\%$\\
\midrule
\textbf{Σύνολο} & & $20$ & $100\%$\\
\bottomrule
\end{tabular}
\end{center}

\elegxos $2+6+6+3+2+1=20$ και $10+30+30+15+10+5=100$.

\paradeigma{Διαβάζω τον πίνακα}
\tixero Ο πίνακας του προηγούμενου παραδείγματος.\par
\tithelo Πόσα παιδιά κάνουν λιγότερο από $10$ λεπτά; Ποιο ποσοστό κάνει τουλάχιστον $15$ λεπτά;\par
\possindeo «Λιγότερο από $10$ λεπτά» σημαίνει οι κλάσεις $0-5$ και $5-10$: $2+6=8$ παιδιά. «Τουλάχιστον $15$» σημαίνει οι κλάσεις $15-20$, $20-25$, $25-30$: $3+2+1=6$ παιδιά, δηλαδή $\frac{6}{20}=0{,}3=30\%$. Ή, με τις σχετικές συχνότητες: $15\%+10\%+5\%=30\%$.\par
\apantisi $8$ παιδιά και $30\%$.\par
\elegxos Τα παιδιά που κάνουν λιγότερο από $10$ λεπτά είναι $40\%$ του συνόλου και τα υπόλοιπα $60\%$.

\paradeigma{Οι κλάσεις μιας μέτρησης}
\tixero Οι μετρήσεις είναι από $12$ έως $47$ και θέλουμε κλάσεις πλάτους $10$.\par
\tithelo Οι κλάσεις και το κέντρο καθεμιάς.\par
\possindeo Ξεκινάμε από το $10$: $10-20$, $20-30$, $30-40$, $40-50$. Το $47$ είναι μέσα στην τελευταία. Το κέντρο κάθε κλάσης είναι το μέσο των άκρων της.\par
\apantisi Τέσσερις κλάσεις με κέντρα $15$, $25$, $35$, $45$.\par
\elegxos Κάθε κέντρο απέχει $5$ (το μισό του πλάτους) από τα άκρα της κλάσης του.

\tmima{Ασκήσεις}

\zoni{Το κατάλαβα}{κλάσεις, κέντρα, συχνότητες}

\ex Σε ποια κλάση από τις $0-10$, $10-20$, $20-30$, $30-40$ ανήκει κάθε τιμή: α) $9$, β) $10$, γ) $19$, δ) $20$, ε) $3$, στ) $30$, ζ) $38$, η) $29{,}9$;

\ex Βρες το πλάτος και το κέντρο των κλάσεων: α) $30-40$, β) $150-160$, γ) $1-1{,}5$, δ) $0-25$.

\ex Συμπλήρωσε τον πίνακα, όπου η μια στήλη έχει μόνο τις συχνότητες $3$, $7$, $6$, $4$ για τις κλάσεις $0-5$, $5-10$, $10-15$, $15-20$: το σύνολο και τις σχετικές συχνότητες.

\ex Η μικρότερη τιμή είναι $7$ και η μεγαλύτερη $28$. Ποιες κλάσεις πλάτους $5$ καλύπτουν όλα τα δεδομένα, αν ξεκινήσουμε από το $5$; Πόσες είναι;

\ex Σε σύνολο $40$ παρατηρήσεων οι συχνότητες πέντε κλάσεων είναι $4$, $10$, $16$, $8$, $2$. Βρες τις σχετικές συχνότητες ως ποσοστά.

\ex Σε σύνολο $50$ παρατηρήσεων οι σχετικές συχνότητες πέντε κλάσεων είναι $8\%$, $20\%$, $36\%$, $24\%$, $12\%$. Βρες τις συχνότητες.

\ex Σωστό ή λάθος; α) Η τιμή $160$ ανήκει στην κλάση $160-165$. β) Το κέντρο της κλάσης $10-20$ είναι $15$. γ) Το άθροισμα των σχετικών συχνοτήτων είναι πάντα $100\%$. δ) Από έναν πίνακα με κλάσεις μπορούμε να βρούμε την ακριβή τιμή κάθε παρατήρησης. ε) Αν μια κλάση έχει συχνότητα $6$ σε σύνολο $24$, η σχετική της συχνότητα είναι $25\%$.

\zoni{Το χρησιμοποιώ}{ομαδοποιώ δικά μου δεδομένα}

\ex \textbf{Το κινητό.} Ο Στέλιος ρώτησε $20$ παιδιά πόσα λεπτά κοιτούν το κινητό το απόγευμα:\par
\noindent\mbox{}\hfill 35 \ 95 \ 15 \ 90 \ 45 \ 85 \ 40 \ 80 \ 45 \ 50 \ 60 \ 70 \ 55 \ 60 \ 65 \ 25 \ 30 \ 55 \ 35 \ 75\hfill\mbox{}\par
Ομαδοποίησε σε $5$ κλάσεις πλάτους $20$ ξεκινώντας από το $0$ και φτιάξε πίνακα συχνοτήτων και σχετικών συχνοτήτων.

\ex \textbf{Οι βαθμοί.} Η κυρία Αλεξάνδρα έβαλε βαθμούς (στα $20$) σε $25$ μαθητές:\par
\noindent\mbox{}\hfill 10 \ 13 \ 15 \ 11 \ 19 \ 8 \ 14 \ 6 \ 9 \ 18 \ 15 \ 9 \ 12 \ 17 \ 12 \ 16 \ 8 \ 17 \ 12 \ 11 \ 13 \ 13 \ 11 \ 14 \ 7\hfill\mbox{}\par
Φτιάξε πίνακα με κλάσεις πλάτους $4$ από το $4$ ως το $20$. Ποιο ποσοστό των μαθητών πήρε τουλάχιστον $12$;

\ex \textbf{Οι πόντοι.} Ο Νίκος σημείωσε πόσους πόντους έβαλε σε $20$ αγώνες μπάσκετ:\par
\noindent\mbox{}\hfill 10 \ 13 \ 11 \ 8 \ 6 \ 8 \ 21 \ 14 \ 19 \ 2 \ 4 \ 12 \ 12 \ 9 \ 7 \ 18 \ 5 \ 15 \ 16 \ 23\hfill\mbox{}\par
Ομαδοποίησε σε κλάσεις πλάτους $5$ από το $0$. Σε πόσους αγώνες έβαλε λιγότερους από $10$ πόντους και ποιο ποσοστό είναι αυτό;

\ex \textbf{Η βιβλιοθήκη.} Σε μια βιβλιοθήκη γειτονιάς γράφτηκαν $80$ παιδιά. Ο πίνακας έχει σβηστεί εν μέρει. Συμπλήρωσέ τον.
\begin{center}
\renewcommand{\arraystretch}{1.2}
\begin{tabular}{ccc}
\toprule
\textbf{Ηλικία (έτη)} & \textbf{Συχνότητα} & \textbf{Σχετική συχνότητα}\\
\midrule
$6-8$ & $12$ & $\ldots$\\
$8-10$ & $\ldots$ & $25\%$\\
$10-12$ & $32$ & $\ldots$\\
$12-14$ & $\ldots$ & $20\%$\\
\midrule
\textbf{Σύνολο} & $80$ & $\ldots$\\
\bottomrule
\end{tabular}
\end{center}

\ex \textbf{Οι σχολικές τσάντες.} Το βάρος της τσάντας $40$ παιδιών, σε κιλά, δίνει συχνότητες $6$, $12$, $14$, $6$, $2$ στις κλάσεις $2-4$, $4-6$, $6-8$, $8-10$, $10-12$. α) Πόσα παιδιά έχουν τσάντα κάτω από $6$ kg και ποιο ποσοστό είναι; β) Ποιο ποσοστό έχει τσάντα τουλάχιστον $8$ kg; γ) Πόσα παιδιά έχουν τσάντα από $4$ ως $8$ kg (χωρίς το $8$);

\zoni{Το σκέφτομαι}{γιατί και πώς}

\ex Ο Στέλιος ομαδοποιεί τα $25$ ύψη σε $12$ κλάσεις πλάτους $2$ cm και η Ρήνη σε $2$ κλάσεις πλάτους $15$ cm. Τι πρόβλημα έχει κάθε επιλογή; Ποιος αριθμός κλάσεων είναι προτιμότερος;

\ex Ο Βασίλης έφτιαξε τον πίνακα για $20$ παρατηρήσεις: κλάσεις $0-10$, $10-20$, $20-30$, $30-40$ με συχνότητες $5$, $8$, $6$, $1$ και σχετικές συχνότητες $25\%$, $40\%$, $30\%$, $10\%$. Βρες το λάθος και διόρθωσέ το.

\ex Βρες $10$ τιμές για τον χρόνο ενός παιδιού στην ουρά του κυλικείου (σε λεπτά, από $0$ έως $9$) ώστε στον πίνακα με κλάσεις $0-3$, $3-6$, $6-9$ οι συχνότητες να βγουν $2$, $5$, $3$. Υπάρχει μόνο μία λύση;

\ex Η κλάση $160-165$ έχει συχνότητα $4$. Γράψε δύο διαφορετικά πιθανά σύνολα ύψων που δίνουν αυτή τη συχνότητα. Τι δεν μπορεί να μας πει ο πίνακας;

\zoni{Για το σπίτι}{μετρώ και ομαδοποιώ}

\ex Μέτρησε το μήκος του πέλματος $10$ ατόμων της οικογένειας ή των φίλων σου, σε εκατοστά, και ομαδοποίησε τις τιμές σε κλάσεις πλάτους $2$ cm. Φτιάξε πίνακα συχνοτήτων και σχετικών συχνοτήτων.

\ex Σημείωσε για μια εβδομάδα πόσα λεπτά διαβάζεις ή μελετάς κάθε μέρα. Ομαδοποίησε τις $7$ τιμές σε $3$ κλάσεις ίσου πλάτους και φτιάξε τον πίνακα.

\paragrafos{Κυκλικά διαγράμματα και ιστογράμματα}{Γ.1.3, «Κυκλικά διαγράμματα – Ιστογράμματα συχνοτήτων κλάσεων ίσου πλάτους με δεδομένο το πλήθος των κλάσεων»}

\begin{anagke}
Η θεία Μαρία έφτιαξε μια πίτσα για τα τρία ανίψια της που είναι $4$, $8$ και $12$ ετών. Θέλει να μοιράσει την πίτσα ανάλογα με την ηλικία τους. Όλοι κοιτούν την πίτσα: τα $4$, $8$ και $12$ κάνουν $24$ μαζί, άρα η μικρότερη πρέπει να πάρει $\frac{4}{24}$ της πίτσας. Αλλά πώς κόβεις ένα κομμάτι που είναι $\frac{4}{24}$ από μια στρογγυλή πίτσα; «Με γωνία από το κέντρο», λέει ο Ηλίας, «το $\frac{4}{24}$ του κύκλου».

Την ίδια ώρα η Ζωή έχει ένα άλλο πρόβλημα. Η παιδίατρος της γειτονιάς σημείωσε την ηλικία των παιδιών που πέρασαν από το ιατρείο μια μέρα και ο πίνακας έχει πάρα πολλές τιμές. «Θέλω ένα σχήμα που να φαίνεται πού μαζεύονται οι ηλικίες», λέει η Ζωή.

Σε αυτή την παράγραφο μαθαίνουμε δύο νέους τρόπους να δείχνουμε δεδομένα με εικόνα: το κυκλικό διάγραμμα, που μοιάζει με πίτσα, και το ιστόγραμμα.
\end{anagke}

\begin{lexi}{κυκλικό, τομέας, επίκεντρη, ιστόγραμμα}{ονόματα για τις εικόνες}
\lexiword{Κυκλικό} είναι ό,τι έχει σχήμα κύκλου. \lexiword{Τομέας} είναι το κομμάτι της πίτσας, από το κέντρο ως την άκρη. Η γωνία του τομέα λέγεται \lexiword{επίκεντρη}, γιατί η κορυφή της είναι στο κέντρο του κύκλου. Το \lexiword{ιστόγραμμα} παίρνει το όνομά του από τον \emph{ιστό}, που σημαίνει «όρθιο στύλο» (όπως ο ιστός ενός πλοίου): τα ορθογώνια του διαγράμματος στέκονται όρθια, σαν στύλοι, το ένα δίπλα στο άλλο.
\end{lexi}

\begin{metafrasi}
Ένα κομμάτι πίτσας. & τομέας κύκλου\\
Η γωνία του κομματιού στο κέντρο. & επίκεντρη γωνία $\hat{\omega}$\\
Ολόκληρος ο κύκλος. & $360^\circ$ και $100\%$\\
Μισή πίτσα. & $180^\circ$ και $50\%$\\
Ορθογώνια κολλητά, ένα για κάθε κλάση. & ιστόγραμμα\\
Πόσο ψηλό είναι το ορθογώνιο. & η συχνότητα της κλάσης\\
\end{metafrasi}

\tmima{Η θεωρία}

\ypotmima{Α}{Κυκλικό διάγραμμα}

\begin{orismos}{κυκλικό διάγραμμα}
Το \textbf{κυκλικό διάγραμμα} είναι ένας κυκλικός δίσκος χωρισμένος σε κυκλικούς τομείς, που παριστάνουν τις διαφορετικές τιμές μιας μεταβλητής.
\end{orismos}

Οι επίκεντρες γωνίες των τομέων είναι ανάλογες με τις συχνότητες των τιμών. Ολόκληρος ο κύκλος αντιστοιχεί στο σύνολο των παρατηρήσεων και έχει γωνία $360^\circ$.

\begin{kanonas}{επίκεντρη γωνία}
$$\frac{\text{συχνότητα της τιμής}}{\text{σύνολο}}=\frac{\hat{\omega}}{360^\circ}\qquad\text{άρα}\qquad \hat{\omega}=\frac{\text{συχνότητα}}{\text{σύνολο}}\cdot 360^\circ$$
Με τη σχετική συχνότητα: $\hat{\omega}=\text{σχετική συχνότητα}\cdot360^\circ$. Το άθροισμα όλων των επίκεντρων γωνιών είναι $360^\circ$.
\end{kanonas}

Το κυκλικό διάγραμμα είναι χρήσιμο για ποσοστά και για λίγες κατηγορίες. Έχει όμως και μειονεκτήματα: είναι δύσκολο να συγκρίνουμε δύο τομείς που διαφέρουν λίγο και δύσκολο να διακρίνουμε πολύ μικρούς τομείς. Όταν οι τιμές είναι πολλές, δεν ταιριάζει.

\ypotmima{Β}{Ιστόγραμμα συχνοτήτων}

Όταν τα δεδομένα έχουν ομαδοποιηθεί σε κλάσεις ίσου πλάτους, τα παριστάνουμε με \textbf{ιστόγραμμα}. 

\begin{orismos}{ιστόγραμμα}
Το \textbf{ιστόγραμμα συχνοτήτων} αποτελείται από συνεχόμενα ορθογώνια που έχουν \textbf{βάση} ίση με το πλάτος κάθε κλάσης και \textbf{ύψος} ίσο με τη συχνότητα της αντίστοιχης κλάσης. Τα άκρα των κλάσεων γράφονται στον οριζόντιο άξονα. Με τον ίδιο τρόπο φτιάχνεται και το \textbf{ιστόγραμμα σχετικών συχνοτήτων}, με ύψη τις σχετικές συχνότητες.
\end{orismos}

Το ιστόγραμμα ταιριάζει σε μεγάλο αριθμό συνεχών δεδομένων που έχουν ομαδοποιηθεί: τα ορθογώνια είναι κολλητά, γιατί το τέλος μιας κλάσης είναι η αρχή της επόμενης, και μας δείχνουν με μια ματιά πού μαζεύονται οι τιμές. Δεν ταιριάζει για κατηγορικά δεδομένα (εκεί χρησιμοποιούμε ραβδόγραμμα) και η εικόνα αλλάζει ανάλογα με τον αριθμό των κλάσεων που θα διαλέξουμε.

\begin{prosoxi}
\begin{itemize}[leftmargin=1.2em,itemsep=0.3ex,topsep=2pt]
\item Στο ραβδόγραμμα οι ράβδοι έχουν κενά, στο ιστόγραμμα τα ορθογώνια ακουμπούν.
\item Στο κυκλικό διάγραμμα οι γωνίες βρίσκονται με τον τύπο $\hat{\omega}=\frac{\text{συχνότητα}}{\text{σύνολο}}\cdot360^\circ$. Μην ξεχνάς το $360^\circ$ ή να διαιρέσεις με το σύνολο.
\item Το ύψος του ορθογωνίου στο ιστόγραμμα είναι η συχνότητα, όχι το πλάτος της κλάσης.
\item Το άθροισμα των γωνιών ενός κυκλικού διαγράμματος πρέπει να βγαίνει $360^\circ$. Αν δεν βγαίνει, έχει γίνει λάθος.
\end{itemize}
\end{prosoxi}

\tmima{Λυμένα παραδείγματα}

\paradeigma{Κατασκευή κυκλικού διαγράμματος}
\tixero Η Φανή ρώτησε $30$ παιδιά τι πίνουν στο διάλειμμα: νερό $12$, χυμό $9$, γάλα $6$, κάτι άλλο $3$.\par
\tithelo Την επίκεντρη γωνία κάθε τομέα και το κυκλικό διάγραμμα.\par
\possindeo Με τον τύπο $\hat{\omega}=\frac{\text{συχνότητα}}{30}\cdot360^\circ$: νερό $\frac{12}{30}\cdot360^\circ=144^\circ$, χυμός $\frac{9}{30}\cdot360^\circ=108^\circ$, γάλα $\frac{6}{30}\cdot360^\circ=72^\circ$, άλλο $\frac{3}{30}\cdot360^\circ=36^\circ$.\par
\apantisi

\begin{center}
\renewcommand{\arraystretch}{1.25}
\begin{tabular}{lccc}
\toprule
\textbf{Ρόφημα} & \textbf{Συχνότητα} & \textbf{Σχετική συχνότητα} & \textbf{Γωνία}\\
\midrule
Νερό & $12$ & $40\%$ & $144^\circ$\\
Χυμός & $9$ & $30\%$ & $108^\circ$\\
Γάλα & $6$ & $20\%$ & $72^\circ$\\
Άλλο & $3$ & $10\%$ & $36^\circ$\\
\midrule
\textbf{Σύνολο} & $30$ & $100\%$ & $360^\circ$\\
\bottomrule
\end{tabular}
\end{center}

\begin{center}
\gaKykliko{1.4}{Νερό/12/$144^\circ$,Χυμός/9/$108^\circ$,Γάλα/6/$72^\circ$,Άλλο/3/$36^\circ$}
\end{center}

\elegxos $144+108+72+36=360$ και $40+30+20+10=100$.

\paradeigma{Διαβάζω κυκλικό διάγραμμα}
\tixero Στο παρακάτω διάγραμμα φαίνονται οι $240$ επισκέπτες μιας σχολικής γιορτής. Οι γωνίες των τομέων είναι: ενήλικες $180^\circ$, παιδιά $90^\circ$, ηλικιωμένοι $60^\circ$, άλλοι $30^\circ$.\par
\tithelo Πόσοι επισκέπτες είναι σε κάθε κατηγορία;\par
\possindeo Από τον τύπο, $\text{συχνότητα}=\frac{\hat{\omega}}{360^\circ}\cdot240$. Ενήλικες: $\frac{180}{360}\cdot240=120$. Παιδιά: $\frac{90}{360}\cdot240=60$. Ηλικιωμένοι: $\frac{60}{360}\cdot240=40$. Άλλοι: $\frac{30}{360}\cdot240=20$.\par
\apantisi $120$, $60$, $40$ και $20$ επισκέπτες.

\begin{center}
\gaKykliko{1.4}{Ενήλικες/180/$180^\circ$,Παιδιά/90/$90^\circ$,Ηλικιωμένοι/60/$60^\circ$,Άλλοι/30/$30^\circ$}
\end{center}

\elegxos $120+60+40+20=240$.

\paradeigma{Κατασκευή ιστογράμματος}
\tixero Ο πίνακας των υψών της Α$3$ από την προηγούμενη παράγραφο: οι κλάσεις $145-150$, $150-155$, $155-160$, $160-165$, $165-170$, $170-175$ έχουν συχνότητες $1$, $6$, $6$, $6$, $4$, $2$.\par
\tithelo Το ιστόγραμμα συχνοτήτων και το ιστόγραμμα σχετικών συχνοτήτων.\par
\possindeo Στον οριζόντιο άξονα βάζουμε τα άκρα των κλάσεων. Φτιάχνουμε ένα ορθογώνιο πάνω από κάθε κλάση, με ύψος ίσο με τη συχνότητά της. Για το ιστόγραμμα σχετικών συχνοτήτων βάζουμε ύψη $4\%$, $24\%$, $24\%$, $24\%$, $16\%$, $8\%$.\par
\apantisi

\begin{center}
\gaIstogramma{145}{5}{6}{1}{1,6,6,6,4,2}{ύψος (cm)}{συχνότητα}\hfill
\gaIstogramma{145}{5}{24}{4}{4,24,24,24,16,8}{ύψος (cm)}{σχετική συχνότητα (\%)}
\end{center}

\elegxos Οι δύο εικόνες έχουν το ίδιο σχήμα. Αλλάζει μόνο η κλίμακα στον κατακόρυφο άξονα.

\paradeigma{Διαβάζω ιστόγραμμα}
\tixero Το ιστόγραμμα δείχνει πόσα λεπτά περίμεναν στον φούρνο $40$ πελάτες.\par
\begin{center}
\gaIstogramma{0}{2}{14}{2}{6,14,12,6,2}{χρόνος αναμονής (λεπτά)}{πελάτες}
\end{center}
\tithelo Πόσοι περίμεναν λιγότερο από $4$ λεπτά; Ποια κλάση έχει τους περισσότερους; Ποιο ποσοστό περίμενε τουλάχιστον $6$ λεπτά;\par
\possindeo Οι κλάσεις έχουν πλάτος $2$. Λιγότερο από $4$ λεπτά σημαίνει οι κλάσεις $0-2$ και $2-4$: $6+14=20$ πελάτες. Ο ψηλότερος στύλος είναι ο $2-4$, με συχνότητα $14$. Τουλάχιστον $6$ λεπτά σημαίνει οι $6-8$ και $8-10$: $6+2=8$ πελάτες, δηλαδή $\frac{8}{40}=20\%$.\par
\apantisi $20$ πελάτες, η κλάση $2-4$, και $20\%$.\par
\elegxos $6+14+12+6+2=40$.

\tmima{Ασκήσεις}

\zoni{Το κατάλαβα}{διαβάζω και υπολογίζω}

\ex Στο κυκλικό διάγραμμα φαίνεται πώς περνούν το Σάββατο τα $40$ παιδιά μιας τάξης.
\begin{center}
\gaKykliko{1.4}{Παρέα/15/15 παιδιά,Οθόνη/10/10 παιδιά,Σπορ/10/10 παιδιά,Βόλτα/5/5 παιδιά}
\end{center}
α) Ποιο ποσοστό των παιδιών κάνει σπορ; β) Ποια είναι η επίκεντρη γωνία του τομέα «Παρέα»; γ) Ποια η γωνία του τομέα «Βόλτα»;

\ex Ρωτήσαμε $40$ μαθητές πώς πηγαίνουν στο σχολείο: με τα πόδια $14$, με λεωφορείο $10$, με αυτοκίνητο $12$, με ποδήλατο $4$. Βρες τις σχετικές συχνότητες και τις επίκεντρες γωνίες των τομέων.

\ex Σε ένα κυκλικό διάγραμμα για τις δραστηριότητες μιας ομάδας παιδιών οι γωνίες είναι: ζωγραφική $90^\circ$, μουσική $120^\circ$, θέατρο $60^\circ$, χορός $60^\circ$, άλλο $30^\circ$. Στη ζωγραφική συμμετέχουν $12$ παιδιά. Πόσα παιδιά έχει η ομάδα και πόσα συμμετέχουν σε κάθε δραστηριότητα;

\ex Στο ιστόγραμμα φαίνονται οι χρόνοι σε δευτερόλεπτα $40$ παιδιών στα $50$ μέτρα ελεύθερο.
\begin{center}
\gaIstogramma{30}{5}{16}{4}{4,10,16,6,4}{χρόνος (s)}{παιδιά}
\end{center}
α) Ποιο είναι το πλάτος κάθε κλάσης; β) Ποια κλάση έχει τη μεγαλύτερη συχνότητα; γ) Πόσα παιδιά έκαναν λιγότερο από $40$ s; δ) Ποιο ποσοστό έκανε τουλάχιστον $45$ s;

\ex Συμπλήρωσε. α) Στο ιστόγραμμα η βάση κάθε ορθογωνίου ισούται με το \ldots\ της κλάσης και το ύψος με τη \ldots\ της κλάσης. β) Ένας τομέας με σχετική συχνότητα $25\%$ έχει επίκεντρη γωνία \ldots. γ) Το άθροισμα όλων των επίκεντρων γωνιών ενός κυκλικού διαγράμματος είναι \ldots. δ) Τα ορθογώνια του ιστογράμματος \ldots\ μεταξύ τους, ενώ οι ράβδοι του ραβδογράμματος όχι.

\ex Ποια αναπαράσταση ταιριάζει καλύτερα: ραβδόγραμμα, κυκλικό διάγραμμα ή ιστόγραμμα; α) Τα ποσοστά που πήρε κάθε υποψήφιος στις εκλογές του μαθητικού συμβουλίου, με τρεις υποψηφίους. β) Τα ύψη $200$ μαθητών, ομαδοποιημένα σε κλάσεις. γ) Πόσα παιδιά προτιμούν κάθε ένα από $8$ διαφορετικά αθλήματα.

\ex Σωστό ή λάθος; α) Ένα ποσοστό $50\%$ αντιστοιχεί σε τομέα $180^\circ$. β) Σε ιστόγραμμα το ύψος κάθε ορθογωνίου είναι ίσο με το πλάτος της κλάσης. γ) Σε κυκλικό διάγραμμα, για σχετική συχνότητα $25\%$ η επίκεντρη γωνία είναι $90^\circ$. δ) Το κυκλικό διάγραμμα ταιριάζει όταν οι κατηγορίες είναι πάρα πολλές. ε) Ένα ιστόγραμμα συχνοτήτων και ένα ιστόγραμμα σχετικών συχνοτήτων των ίδιων δεδομένων έχουν το ίδιο σχήμα.

\zoni{Το χρησιμοποιώ}{φτιάχνω διαγράμματα}

\ex \textbf{Το φύλλο.} Η θεία Μαρία μέτρησε το μήκος $20$ φύλλων από το δέντρο της. Οι κλάσεις $4-6$, $6-8$, $8-10$, $10-12$, $12-14$ (σε cm) έχουν συχνότητες $2$, $5$, $8$, $4$, $1$. Φτιάξε το ιστόγραμμα συχνοτήτων.

\ex \textbf{Το βιβλιοπωλείο.} Σε $60$ αγορές ενός βιβλιοπωλείου: τετράδια $24$, στυλό $15$, μπλοκ $12$, άλλα είδη $9$. Βρες τις επίκεντρες γωνίες και φτιάξε το κυκλικό διάγραμμα.

\ex \textbf{Οι τσάντες.} Στην παράγραφο Γ1.2 τα $40$ παιδιά είχαν συχνότητες τσάντας $6$, $12$, $14$, $6$, $2$ στις κλάσεις $2-4$, $4-6$, $6-8$, $8-10$, $10-12$. Βρες τις σχετικές συχνότητες και φτιάξε το ιστόγραμμα σχετικών συχνοτήτων.

\ex \textbf{Οι προτιμήσεις.} Στο κυκλικό διάγραμμα φαίνονται οι προτιμήσεις $36$ παιδιών για τρία αθλήματα και μία κατηγορία «άλλο».
\begin{center}
\gaKykliko{1.4}{Ποδόσφαιρο/150/$150^\circ$,Μπάσκετ/90/$90^\circ$,Κολύμβηση/70/$70^\circ$,Άλλο/50/$50^\circ$}
\end{center}
Πόσα παιδιά προτιμούν κάθε άθλημα; Ποιο κλάσμα των παιδιών προτιμά το ποδόσφαιρο;

\ex \textbf{Το μουσείο.} Το ιστόγραμμα δείχνει τις ηλικίες $60$ επισκεπτών ενός μουσείου.
\begin{center}
\gaIstogramma{0}{10}{16}{4}{6,15,12,9,12,6}{ηλικία (έτη)}{επισκέπτες}
\end{center}
α) Πόσοι επισκέπτες είναι κάτω από $20$ ετών και ποιο ποσοστό είναι αυτό; β) Ποιο ποσοστό είναι από $40$ ετών και πάνω; γ) Σε ποια κλάση ανήκει η πιο συνηθισμένη ηλικία;

\zoni{Το σκέφτομαι}{γιατί και πότε}

\ex Η Αγγελική έφτιαξε κυκλικό διάγραμμα για το αγαπημένο χρώμα των συμμαθητών της, με $9$ διαφορετικά χρώματα και μερικά από αυτά με πολύ μικρούς τομείς. Τι δυσκολίες θα έχει όποιος το διαβάσει; Ποιο άλλο διάγραμμα θα πρότεινες;

\ex Εξήγησε γιατί το $1\%$ αντιστοιχεί σε $3{,}6^\circ$ στο κυκλικό διάγραμμα. Βρες την επίκεντρη γωνία για ποσοστά $5\%$, $12\%$ και $37{,}5\%$.

\ex Στο ιστόγραμμα της Α$3$ για τα ύψη, ο Ηλίας ενώνει ανά δύο τις γειτονικές κλάσεις και φτιάχνει τρεις κλάσεις πλάτους $10$: $145-155$, $155-165$, $165-175$. Βρες τις νέες συχνότητες και σχεδίασε το νέο ιστόγραμμα. Τι αλλάζει στην εικόνα και τι μένει ίδιο;

\ex Ο Βασίλης υπολόγισε τις γωνίες για $24$ παιδιά: $8$ παιδιά $120^\circ$, $6$ παιδιά $90^\circ$, $7$ παιδιά $110^\circ$, $3$ παιδιά $45^\circ$. Βρες ποια γωνία είναι λάθος και διόρθωσέ την.

\zoni{Για το σπίτι}{η μέρα μου σε κύκλο}

\ex Χώρισε σε κυκλικό διάγραμμα τις $24$ ώρες της μέρας σου: ύπνος, σχολείο, φαγητό και μετακινήσεις, διάβασμα, ελεύθερος χρόνος. Βρες πρώτα πόσες μοίρες αντιστοιχούν σε μία ώρα.

\ex Βρες σε εφημερίδα ή ιστοσελίδα ένα κυκλικό διάγραμμα ή ένα ιστόγραμμα και γράψε τρία συμπεράσματα που προκύπτουν από αυτό.

\epanalipsi{Γ1}

\tmima{Λεξικό του κεφαλαίου}

\renewcommand{\arraystretch}{1.5}
\begin{longtable}{@{}>{\raggedright\arraybackslash}p{0.28\linewidth} >{\raggedright\arraybackslash}p{0.62\linewidth}@{}}
\toprule
\textbf{Λέξη} & \textbf{Με μία πρόταση}\\
\midrule
\endhead
πληθυσμός, μεταβλητή & Αυτό που μελετάμε και το χαρακτηριστικό του που μας ενδιαφέρει.\\
απογραφή, δείγμα & Εξετάζω όλο τον πληθυσμό· εξετάζω ένα μέρος του.\\
αντιπροσωπευτικό δείγμα & Μικρογραφία του πληθυσμού, που δίνει αξιόπιστα συμπεράσματα.\\
κατηγορικά, ποσοτικά δεδομένα & Δεν είναι αριθμοί· είναι αριθμοί.\\
διακριτά, συνεχή & Μεμονωμένες τιμές (πόσα)· οποιαδήποτε τιμή σε διάστημα (πόσο).\\
συχνότητα, σχετική συχνότητα & Πόσες φορές εμφανίζεται μια τιμή· συχνότητα διά το σύνολο, ως ποσοστό.\\
κλάση, πλάτος, κέντρο & Ομάδα τιμών· διαφορά των άκρων της· μέσο των άκρων της.\\
ραβδόγραμμα & Ράβδοι με κενά, ύψος ίσο με τη συχνότητα. Για κατηγορικά και διακριτά.\\
κυκλικό διάγραμμα & Τομείς με $\hat{\omega}=\frac{\text{συχνότητα}}{\text{σύνολο}}\cdot360^\circ$.\\
ιστόγραμμα & Κολλητά ορθογώνια με βάση το πλάτος της κλάσης και ύψος τη συχνότητά της.\\
\bottomrule
\end{longtable}

\tmima{Ασκήσεις επανάληψης}

\ex Κατηγορικά ή ποσοτικά δεδομένα; Στα ποσοτικά πες αν είναι διακριτά ή συνεχή. α) Ο αριθμός των κατοικίδιων μιας οικογένειας. β) Το χρώμα της μπλούζας. γ) Ο χρόνος κολύμβησης στα $100$ μέτρα. δ) Το ύψος ενός δέντρου. ε) Το πλήθος των μαθητών ενός τμήματος. στ) Η αγαπημένη εφαρμογή στο κινητό.

\ex Η Αγγελική θέλει να μάθει τι είδους μουσική ακούν οι $300$ μαθητές του σχολείου της και ρωτά $40$ παιδιά που επιλέγει με κλήρωση από όλες τις τάξεις. Ποιος είναι ο πληθυσμός, ποιο το δείγμα, ποια η μεταβλητή και τι είδους είναι τα δεδομένα;

\ex Ο Νίκος θέλει να μάθει πόσα χρήματα ξοδεύουν οι μαθητές του σχολείου στο κυλικείο και ρωτά τα παιδιά που περιμένουν στην ουρά του κυλικείου. Είναι το δείγμα του αντιπροσωπευτικό; Δικαιολόγησε.

\ex Η Ζωή ρώτησε $20$ συμμαθητές ποιο είναι το αγαπημένο τους φρούτο: Μήλο (Μ), Βερίκοκο (Β), Πορτοκάλι (Π), Φράουλα (Φ), Κεράσι (Κ).\par
\noindent Β Π Κ Μ Π Π Β Φ Φ Μ Μ Μ Β Β Μ Κ Β Φ Π Μ\par
Φτιάξε πίνακα συχνοτήτων με σχετικές συχνότητες και ραβδόγραμμα.

\ex Το μήκος άλματος (σε cm) $20$ παιδιών στο μάθημα της γυμναστικής:\par
\noindent\mbox{}\hfill 151 \ 160 \ 138 \ 150 \ 172 \ 158 \ 142 \ 176 \ 167 \ 140 \ 170 \ 161 \ 181 \ 147 \ 188 \ 165 \ 145 \ 132 \ 155 \ 153\hfill\mbox{}\par
Ομαδοποίησε τα δεδομένα σε κλάσεις πλάτους $10$ ξεκινώντας από το $130$ και φτιάξε πίνακα συχνοτήτων και σχετικών συχνοτήτων.

\ex Για τις κλάσεις της προηγούμενης άσκησης βρες το πλάτος και το κέντρο καθεμιάς.

\ex Συμπλήρωσε τον πίνακα για $40$ παρατηρήσεις.
\begin{center}
\renewcommand{\arraystretch}{1.2}
\begin{tabular}{ccc}
\toprule
\textbf{Κλάσεις} & \textbf{Συχνότητα} & \textbf{Σχετική συχνότητα}\\
\midrule
$0-10$ & $4$ & $\ldots$\\
$10-20$ & $\ldots$ & $30\%$\\
$20-30$ & $14$ & $\ldots$\\
$30-40$ & $\ldots$ & $20\%$\\
$40-50$ & $2$ & $\ldots$\\
\midrule
\textbf{Σύνολο} & $40$ & $\ldots$\\
\bottomrule
\end{tabular}
\end{center}

\ex Με τα δεδομένα του πίνακα της άσκησης $\text{Ε}\Gamma1.5$ βρες πόσα παιδιά πήδηξαν τουλάχιστον $160$ cm και ποιο ποσοστό πήδηξε λιγότερο από $150$ cm.

\ex Ρωτήσαμε $24$ παιδιά ποιο παιχνίδι προτιμούν στο διάλειμμα: κρυφτό $8$, μπάλα $6$, σχοινάκι $6$, κυνηγητό $4$. Βρες τις επίκεντρες γωνίες και φτιάξε το κυκλικό διάγραμμα.

\ex Σε κυκλικό διάγραμμα για $72$ μαθητές οι γωνίες των τομέων είναι: βόλεϊ $100^\circ$, μπάσκετ $80^\circ$, τένις $60^\circ$, ποδόσφαιρο $120^\circ$. Πόσοι μαθητές προτιμούν κάθε άθλημα;

\ex Οι $25$ βαθμοί της κυρίας Αλεξάνδρας στην παράγραφο Γ1.2 έδωσαν συχνότητες $2$, $8$, $10$, $5$ στις κλάσεις $4-8$, $8-12$, $12-16$, $16-20$. Φτιάξε το ιστόγραμμα συχνοτήτων και το ιστόγραμμα σχετικών συχνοτήτων.

\ex Στο ιστόγραμμα φαίνονται τα λεπτά μελέτης την ημέρα $50$ παιδιών.
\begin{center}
\gaIstogramma{0}{20}{20}{5}{5,12,18,10,5}{μελέτη (λεπτά)}{παιδιά}
\end{center}
α) Πόσα παιδιά μελετούν λιγότερο από $40$ λεπτά και ποιο ποσοστό είναι αυτό; β) Πόσα μελετούν τουλάχιστον $60$ λεπτά; γ) Ποια κλάση έχει τη μεγαλύτερη συχνότητα;

\ex Διάλεξε το κατάλληλο διάγραμμα (ραβδόγραμμα, κυκλικό ή ιστόγραμμα) και εξήγησε: α) τα βάρη $150$ μαθητών, ομαδοποιημένα σε κλάσεις, β) τα μερίδια τεσσάρων εφαρμογών σε ένα σύνολο $100\%$, γ) πόσα παιδιά προτιμούν καθένα από $7$ είδη παγωτού.

\ex Η Πολυξένη βρήκε για τέσσερις κατηγορίες σχετικές συχνότητες $20\%$, $30\%$, $25\%$ και $15\%$ και έφτιαξε κυκλικό διάγραμμα. Τι δεν πάει καλά; Αν το λάθος είναι στην τρίτη κατηγορία, ποιες είναι οι σωστές γωνίες των τομέων;

\ex Σωστό ή λάθος; α) Τα συνεχή δεδομένα παίρνουν μεμονωμένες τιμές. β) Ο αριθμός των αδελφών ενός παιδιού είναι ποσοτικό διακριτό δεδομένο. γ) Η απογραφή εξετάζει όλα τα στοιχεία του πληθυσμού. δ) Ένα αντιπροσωπευτικό δείγμα δίνει αξιόπιστα συμπεράσματα. ε) Το χρώμα ματιών είναι ποσοτικό. στ) Η συχνότητα μιας τιμής είναι πάντα μικρότερη από το σύνολο ή ίση με αυτό. ζ) Το άθροισμα των σχετικών συχνοτήτων είναι πάντα $100\%$. η) Το δεξί άκρο μιας κλάσης ανήκει στην ίδια κλάση.

\ex Σωστό ή λάθος; α) Στο κυκλικό διάγραμμα ένα ποσοστό $25\%$ αντιστοιχεί σε γωνία $90^\circ$. β) Στο ιστόγραμμα τα ορθογώνια έχουν κενά. γ) Σε ένα ιστόγραμμα με κλάσεις ίσου πλάτους, το ύψος κάθε ορθογωνίου είναι η συχνότητα της κλάσης. δ) Το κέντρο της κλάσης $10-20$ είναι $15$. ε) Το κυκλικό διάγραμμα ταιριάζει σε πολλές κατηγορίες. στ) Στο ραβδόγραμμα οι ράβδοι ακουμπούν. ζ) Το ιστόγραμμα σχετικών συχνοτήτων έχει το ίδιο σχήμα με το ιστόγραμμα συχνοτήτων των ίδιων δεδομένων. η) Οι επίκεντρες γωνίες ενός κυκλικού διαγράμματος έχουν άθροισμα $180^\circ$.

\Needspace{17\baselineskip}\tmima{Ο χάρτης της επανάληψης}

Μετά τις λύσεις, δες σε ποιες ασκήσεις δυσκολεύτηκες. Ο πίνακας σου λέει ποια παράγραφο να ξαναδιαβάσεις.

\renewcommand{\arraystretch}{1.5}
\begin{tabular*}{\linewidth}{@{\extracolsep{\fill}} >{\raggedright\arraybackslash}p{0.38\linewidth} >{\raggedright\arraybackslash}p{0.52\linewidth}}
\toprule
\textbf{Αν δυσκολεύτηκες στις} & \textbf{ξαναδιάβασε τις παραγράφους}\\
\midrule
ΕΓ1.1, ΕΓ1.2, ΕΓ1.3, ΕΓ1.4, ΕΓ1.15 & 1.1 (πληθυσμός, δεδομένα, ραβδόγραμμα)\\
ΕΓ1.5, ΕΓ1.6, ΕΓ1.7, ΕΓ1.8 & 1.2 (κλάσεις και πίνακες)\\
ΕΓ1.9, ΕΓ1.10, ΕΓ1.11, ΕΓ1.12, ΕΓ1.13, ΕΓ1.14, ΕΓ1.16 & 1.3 (κυκλικά διαγράμματα και ιστογράμματα)\\
\bottomrule
\end{tabular*}

